%==============================================================================
\chapter{Délais et livrables}
\label{chap:delais}
%==============================================================================

\section{Livrables attendus}
\begin{itemize}
  \item Application fonctionnelle (code source serveur Java + front-end).
  \item Plan de tests et résultats d'exécution.
  \item Dossier de conception (diagrammes UML).
  \item Rapport de projet, poster et présentation orale.
\end{itemize}

\section{Critères de recette}
La recette est prononcée lorsque \textbf{l'ensemble} des critères suivants est
satisfait et constaté contradictoirement. Chaque critère est vérifiable :
un critère qu'on ne sait pas mesurer n'est pas un critère de recette.

\begin{xltabular}{\textwidth}{>{\bfseries}p{3.2cm} X p{3.4cm}}
\toprule
\textbf{Domaine} & \textbf{Critère} & \textbf{Moyen de preuve} \\
\midrule
\endhead
Fonctionnel & Les cas d'utilisation du chapitre~\ref{chap:cas-utilisation} sont réalisés et démontrés sur données réelles. & Procès-verbal de recette \\
\addlinespace
Qualité de code & Couverture de tests de la couche métier $\geq$ 70~\%, aucune vulnérabilité critique ou haute non traitée. & Rapport JaCoCo, scan de dépendances \\
\addlinespace
Sécurité & Test d'intrusion externe réalisé ; aucune vulnérabilité critique ou haute ouverte à la livraison. & Rapport de pentest et plan de remédiation \\
\addlinespace
Protection des données & Coordonnées de géolocalisation inaccessibles hors rôle autorisé, absentes des exports par défaut et des journaux. & Tests d'autorisation, revue des journaux \\
\addlinespace
Performance & p95 $<$ 500~ms et taux d'erreur $<$ 1~\% sous la charge nominale définie à l'annexe~\ref{chap:exploitation}. & Rapport de tests de charge k6 \\
\addlinespace
Disponibilité & Les objectifs de service (SLO), de reprise (RTO) et de perte maximale (RPO) de l'annexe~\ref{chap:exploitation} sont tenus. & Supervision, test de restauration \\
\addlinespace
Sauvegarde & Une restauration complète a été réalisée avec succès en conditions réelles. & Compte rendu de test de restauration \\
\addlinespace
Accessibilité et ergonomie & Interface utilisable au clavier, contrastes conformes à la charte, parcours validé par un utilisateur métier. & Audit d'accessibilité, recette utilisateur \\
\addlinespace
Internationalisation & Les trois langues (FR, EN, AR) sont complètes, le rendu RTL arabe est vérifié sur tous les écrans. & Revue linguistique \\
\addlinespace
Exploitabilité & Runbook livré, transfert de compétences effectué, système redéployable depuis le dépôt. & Séance de transfert, redéploiement à blanc \\
\addlinespace
Réversibilité & Données exportables dans un format ouvert et documenté ; code source et documentation remis. & Export de démonstration \\
\bottomrule
\end{xltabular}

\begin{encadre}
\textbf{Réserves.} Une réserve mineure n'empêche pas la recette mais fait
l'objet d'un délai de levée convenu. Une réserve \textbf{bloquante} --- toute
vulnérabilité critique, tout échec de restauration, toute fuite de données entre
exploitations --- suspend la recette jusqu'à correction et nouvelle vérification.
\end{encadre}

\section{Diagrammes UML à produire}
Diagrammes de : cas d'utilisation, classes, objets, séquence, activité, machine
à états, interaction / collaboration, package, composant et déploiement. Le
contenu attendu de chacun est détaillé au chapitre \ref{chap:conception} portant
sur la conception.

\section{Planning prévisionnel}
Le planning est exprimé en semaines relatives, avec les livrables intermédiaires
associés à chaque phase.
\begin{xltabular}{\textwidth}{>{\bfseries}p{4.5cm} C{2.2cm} X}
\toprule
\textbf{Phase} & \textbf{Semaines} & \textbf{Livrable} \\
\midrule
\endhead
Analyse et cahier des charges & S1 à S2 & Cahier des charges validé \\
Conception (UML) & S3 à S4 & Dossier de conception \\
Développement du noyau CRUD & S5 à S7 & Modèle et gestion des entités \\
Tableaux de bord et alertes & S8 à S9 & Calendrier, alertes et synthèses \\
Services web, sécurité, i18n & S10 à S11 & OAuth, API, X.509 et SSL, XML \\
Tests et validation & S12 & Plan de tests exécuté \\
Documentation et présentation & S13 & Rapport, poster et soutenance \\
\bottomrule
\end{xltabular}

